\documentclass[11pt]{article}
\usepackage{preamble}

\newcommand{\IconCheck}{[CHECK]}
\newcommand{\IconMic}{[MIC]}
\newcommand{\IconClipboard}{[CLIPBOARD]}

\newcommand{\phasetag}[1]{%
  {\small\itshape Tag: #1\par}
}

\newcommand{\phaseitems}[1]{%
  \begin{itemize}[leftmargin=1.5em,itemsep=2pt,topsep=2pt]
  #1
  \end{itemize}%
}

\newcommand{\teacheranswerbox}[2]{%
  \begin{callout}{#1}{calloutgreen}
  #2
  \end{callout}
}

\begin{document}

\packetheading{Lesson 4-1 · Period 2 · Teacher Edition}{Inverse Variation Applications + DOK-3 Performance Task}

\sectionbanner{OBJECTIVES}
{\small
\renewcommand{\arraystretch}{1.25}
\noindent\begin{tabular}{|p{1.8in}|p{4.8in}|}
\hline
\textbf{Math Objective} & Apply inverse variation to real-world contexts (wavelength, gas pressure, travel distance), and distinguish DIRECT variation (\(y = kx\)) from INVERSE variation (\(xy = k\)) inside a single multi-part problem. \\ \hline
\textbf{Language Objective} & Students will use the frame “This is \_\_\_ variation because as \_\_\_ \_\_\_, \_\_\_ \_\_\_ proportionally.” to justify whether a relationship is direct or inverse. \\ \hline
\textbf{Essential Understanding} & Not every covariation is inverse variation. Direct variation has a constant RATIO (\(y/x = k\)); inverse variation has a constant PRODUCT (\(xy = k\)). A single scenario can contain BOTH types across different variable pairs. \\ \hline
\end{tabular}
}

\sectionbanner{TOPIC GOALS / ESSENTIAL QUESTION / MATERIALS}
{\small
\renewcommand{\arraystretch}{1.25}
\noindent\begin{tabular}{|p{1.8in}|p{4.8in}|}
\hline
\textbf{Topic Goals} & Students model 4 contextual inverse-variation problems, then tackle a Performance Task (Ramón's road trip) that embeds direct and inverse variation in one scenario. Every item is Savvas-sourced. \\ \hline
\textbf{Essential Question} & How does inverse variation model real-world relationships, and how is it different from direct variation? \\ \hline
\textbf{Materials} & Pencils; student packet; calculator. No Chromebooks required. \\ \hline
\end{tabular}
}

\begin{callout}{\IconPin\ PRE-FLIGHT — P2 IS THE DOK-3 SPINE DAY}{calloutpink}
Today's conceptual core is Practice \#26 (Ramón). Part A is direct variation, Part B is inverse — same scenario, different variable pairs. That contrast IS the lesson's deepest content.\\
Pace: Try-It 3 (5 min) \(\rightarrow\) \#12 (4 min) \(\rightarrow\) \#17 (5 min) \(\rightarrow\) \#23 (7 min) \(\rightarrow\) \#26 (12--15 min). \(\sim\)33--36 min total Explore.\\
If running tight: cut \#12 (the HOT inverse-square is tangent, not central). Never cut \#26.
\end{callout}

\phasetag{bridge from P1}
\frameworkphaseheader{Do Now}{1}{5}{%
\phaseitems{
  \item Hand out at door.
  \item Circulate 3 min.
  \item Cold-call 1 student for \(xy\) computation.
}%
}{%
\phaseitems{
  \item Compute \(xy\) for each column.
  \item State \(k\) if constant; explain if not.
}%
}{%
\phaseitems{
  \item What do you multiply to check inverse variation?
  \item Is \(xy\) constant here? What does that tell you?
}%
}{Solo teacher: circulate silently. No hints.}

\bankitem{Do Now (4-1-savvas-q15)}{%
15. Do the tables of values represent inverse variations? Explain. See Example 1. Table - \(x\): \(1, 2, 3, 4, 5, 6\); \(y\): \(60, 30, 20, 15, 12, 10\).
\begin{center}
\renewcommand{\arraystretch}{1.2}
\begin{tabular}{|c|c|c|c|c|c|c|}
\hline
\(x\) & 1 & 2 & 3 & 4 & 5 & 6 \\ \hline
\(y\) & 60 & 30 & 20 & 15 & 12 & 10 \\ \hline
\end{tabular}
\end{center}
}{}

\teacheranswerbox{\IconCheck\ ANSWER KEY}{%
\textbf{Answer key (from TE):} Yes; as the \(x\)-values increase, the \(y\)-values decrease. Also the product of each pair of \(x\)- and \(y\)-values is \(60\) (\(k = 60\), so \(y = 60/x\)).
}

\phasetag{teacher models Example 3}
\frameworkphaseheader{Launch}{2}{12}{%
\phaseitems{
  \item Project Example 3 (bouzouki). Think aloud: extract the two quantities (\(s =\) string length, \(f =\) frequency). Compute \(k = sf\).
  \item “Now I can find \(f\) for any \(s\). Plug in the new length, solve.”
  \item Flag the Application Rules card — print it in student minds.
  \item 30-sec pause: “What's \(k\) for YOUR Do Now table?” (partner check)
  \item Do NOT solve Try-It 3 or Practice items — release at minute 17.
}%
}{%
\phaseitems{
  \item Watch the bouzouki example silently.
  \item During pause: state \(k\) from Do Now to partner.
  \item Copy Application Rules card into margin.
  \item Note the DIRECT vs INVERSE warning: this comes back in \#26.
}%
}{%
\phaseitems{
  \item What two quantities are related? What's constant?
  \item “Every radio wave with frequency \(f\) has wavelength \(w\) such that \(wf =\) constant.” Why is that inverse?
  \item If frequency doubles, what happens to wavelength? Why?
}%
}{Project. Think aloud. Release promptly at minute 17.}

\begin{callout}{\IconMic\ TEACHER SCRIPT}{calloutblue}
1. “Watch me set up a real-world inverse variation. Bouzouki string: length \(s\) varies inversely with frequency \(f\).”\\
2. “Step 1: Identify the quantities and find \(k\). \(k = sf = 26 \times 329.63 = 8{,}570.38\).”\\
3. “Step 2: Write the equation. \(s = k/f\), so \(s = 8570.38/f\).”\\
4. “Step 3: Plug in the new value. For \(s = 13\): \(f = 8570.38/13 = 659.26\) cycles/sec.”\\
5. “Tonight's Performance Task has a twist — ONE scenario contains BOTH direct and inverse variation. Be ready to tell them apart.”\\
6. Release to Explore.
\end{callout}

\phasetag{TPS + DOK-3 driver}
\frameworkphaseheader{Explore}{2--3}{33}{%
\phaseitems{
  \item 5 laps across 33 min.
  \item Lap 1 (5 min): Try-It 3 — ice cube. Press pairs to identify \(k =\) Temp \(\times\) Time.
  \item Lap 2 (4 min): \#12 — inverse square. This is a THINKING item; let pairs struggle.
  \item Lap 3 (5 min): \#17 — graph reading + context. Confirm \(k = 300{,}000\).
  \item Lap 4 (7 min): \#23 — Boyle's Law. Check that pairs compute volleyball's \(k = 1350\).
  \item Lap 5 (12 min): \IconStar\ \#26 Ramón. PART A is direct; PART B is inverse. Press hard on that contrast.
  \item Do NOT give \#26 answers. Pairs present Part A/B to each other; one pair presents at Share.
}%
}{%
\phaseitems{
  \item 5 items in order: Try-It 3 \(\rightarrow\) \#12 \(\rightarrow\) \#17 \(\rightarrow\) \#23 \(\rightarrow\) \#26.
  \item For \#26: BEFORE solving, identify which part is direct and which is inverse. Write it down.
  \item Justify with sentence frame.
}%
}{%
\phaseitems{
  \item What's \(k\) in this problem? What units?
  \item \#12: \(y = k/x^2\). If \(x \rightarrow 4x\), what happens to \(y\)? (Divide by 16.)
  \item \#17: from the graph, what's \(k\)? How do you read it?
  \item \#23: Whose \(k\) do we use — the volleyball's, or the basketball's? (Volleyball's; the basketball only gives the target PRESSURE.)
  \item \#26 Part A: is distance/time constant? \(\rightarrow\) direct variation.
  \item \#26 Part B: is gas \(\times\) time constant? \(\rightarrow\) inverse variation.
}%
}{Solo teacher: 5 laps. Press DIRECT vs INVERSE distinction on \#26.}

\bankitem{Try It 3 (4-1-tryit-3)}{%
Try It 3. The amount of time it takes for an ice cube to melt varies inversely with the air temperature in degrees Celsius. At \(20^\circ\)C, the ice will melt in 20 minutes. How long will it take the ice to melt if the temperature is \(30^\circ\)C?
}{}

\teacheranswerbox{\IconCheck\ ANSWER / ANTICIPATED TALK}{%
\textbf{Answer key (from TE):} \(\approx 13.3\) minutes. Reasoning: \(k = 20\cdot 20 = 400\); \(t = 400/30 \approx 13.3\).
}

\bankitem{Practice \#12 (4-1-savvas-q12)}{%
12. Higher Order Thinking: Suppose \(y\) varies inversely as the square of \(x\). If \(x\) is multiplied by 4, what is true for the value of \(y\)?
}{}

\teacheranswerbox{\IconCheck\ ANSWER / ANTICIPATED TALK}{%
\textbf{Answer key (from TE):} Divide by 16. Reasoning: \(y = k/x^2\). When \(x \rightarrow 4x\), new \(y = k/(4x)^2 = k/(16x^2) = (1/16)y_{\text{original}}\).
}

\begin{callout}{\IconSpeak\ CHAIN 1 CALLBACK --- Wavelength/Frequency (first callback)}{calloutpeach}
\textit{When to say:} Before students start Practice \#17.

\smallskip
``Remember Storage Box? We had a volume constraint, we wanted a dimension. Here we have a \textbf{product constraint} --- wavelength times frequency equals the speed of light --- and we still want to pull out one side given the other. Different equation, same job.''

\smallskip
\textit{Secondary plant:} this is where inverse variation becomes the thematic bridge between polynomial modeling (Topic 3) and the reciprocal/log families (Topics 4--6). Name it.
\end{callout}

\bankitem{Practice \#17 (4-1-savvas-q17)}{%
17. The wavelength, \(w\), of a radio wave varies inversely to its frequency, \(f\), as shown in the graph. The graph shows \(w\) (m) on the y-axis from 0 to 1200, and \(f\) (kilohertz) on the x-axis from 0 to 10000, with the curve \(w = k/f\) passing through data points including (1000, 300), (2000, 150), (4000, 75), (6000, 50), (8000, 37.5), (10000, 30). A radio wave with a frequency of 1,000 kilohertz has a length of 300 m. What is the frequency when the wavelength is 375 m?
\begin{center}
\begin{tikzpicture}
\begin{axis}[
  width=0.9\linewidth,
  height=2.6in,
  xmin=0,
  xmax=10000,
  ymin=0,
  ymax=1200,
  xlabel={Frequency (kilohertz)},
  ylabel={Wavelength (m)},
  grid=major,
  minor tick num=1,
  tick label style={font=\small},
  label style={font=\small},
  every axis plot/.append style={line width=1pt}
]
\addplot[tealheader, smooth, domain=250:10000, samples=250] {300000/x};
\addplot[only marks, mark=*, mark size=2.2pt, color=dayblue] coordinates {
  (1000,300)
  (2000,150)
  (4000,75)
  (6000,50)
  (8000,37.5)
  (10000,30)
};
\node[anchor=south east, font=\small] at (axis cs:9800, 1040) {\(w = \dfrac{k}{f}\)};
\end{axis}
\end{tikzpicture}
\end{center}
}{}

\teacheranswerbox{\IconCheck\ ANSWER / ANTICIPATED TALK}{%
\textbf{Answer key (from TE):} 800 kHz. Transcription and graph description verified against Gemini 3.1 Pro (LaTeX output with \(w = 300000/f\) curve and data points).
}

\bankitem{Practice \#23 (4-1-savvas-q23)}{%
23. Model With Mathematics: Boyle's Law states that the pressure exerted by fixed quantity of a gas, \(P\), varies inversely with the volume the gas occupies, \(V\), assuming constant temperature. The volume and air pressure of a volleyball are 300 in\(^{3}\) and 4.5 psi. The volume and air pressure of a basketball are 415 in\(^{3}\) and 8 psi. How much smaller would the volleyball have to be to be equal the air pressure of the basketball?
}{}

\teacheranswerbox{\IconCheck\ ANSWER / ANTICIPATED TALK}{%
\textbf{Answer key (from TE):} 131.25 in\(^3\). Volleyball \(k_V = 300 \times 4.5 = 1350\). At \(P = 8\) psi, \(V = 1350/8 = 168.75\) in\(^3\). Difference = \(300 - 168.75 = 131.25\) in\(^3\). (Note: the basketball's own \(k\) is irrelevant to the answer - we only use basketball's \(P = 8\) psi as the target for the volleyball.)
}

\bankitem{Practice \#26 \IconStar\ DOK-3 DRIVER (4-1-savvas-q26)}{%
26. Performance Task. Suppose Ramón takes a road trip. He starts from his home in the suburbs of Cleveland, OH and travels to Pittsburgh, PA to visit his aunt and uncle.

Part A: The distance Ramón drives from Cleveland to Pittsburgh is 133 miles. The trip takes him 2 hours. The distance \(d\) in miles that Ramón drives varies directly with the amount of time \(t\) in hours he spends driving. Write the equation of the direct variation. Use the given relationship and the equation to find the number of miles Ramón would travel if he continues on for 5 more hours.

Part B: The amount of gas in Ramón's car is 9 gal after he drives for 2 h. The amount of gas \(g\) in gallons in his tank varies inversely with the amount of time \(t\) in hours he spends driving. Write the equation of the inverse variation. Use the given relationship and the equation to find the number of gallons in Ramón's tank after 5 more hours of driving.

\begin{center}
\begin{tikzpicture}[>=stealth, line width=0.9pt, font=\small]
\node[draw, rounded corners, fill=frameshade, minimum width=1.7in, minimum height=0.58in] (cle) at (0,0) {\textbf{Cleveland, OH}};
\node[draw, rounded corners, fill=frameshade, minimum width=1.7in, minimum height=0.58in] (pit) at (5.0,0) {\textbf{Pittsburgh, PA}};
\draw[->, tealheader] (cle.east) -- node[above, align=center] {133 miles\\2 hours} (pit.west);
\node[draw, rounded corners, fill=calloutyellow, minimum width=2.1in, minimum height=0.6in] at (2.5,-1.0) {After 2 h: 9 gal in tank};
\node[draw, rounded corners, fill=calloutpeach, minimum width=2.3in, minimum height=0.6in] at (2.5,-2.0) {Continue for 5 more hours};
\end{tikzpicture}
\end{center}
}{}

\teacheranswerbox{\IconCheck\ ANSWER / ANTICIPATED TALK}{%
Transcription verified against Gemini 3.1 Pro (clarified “5 more hours” vs my initial “3.5” misread). \textbf{Answer key (from TE):} Part A \(d = 66.5t\); 465.5 mi (total \(t = 7\) hours since \(2 + 5\)). Part B \(g = 18/t\); approximately 2.57 gal (at \(t = 7\)). STRONG DOK-3 driver candidate - the direct-vs-inverse contrast is pedagogically richer than the pure-inverse items \#20/\#22/\#25. Recommend this as the lesson's single-DOK-3 spine.
}

\phasetag{academic conversation + Wed bridge}
\frameworkphaseheader{Share / Summary}{2}{5}{%
\phaseitems{
  \item Cold-call 1 pair to present Ramón's Part A (direct) + Part B (inverse) using sentence frame.
  \item Class signals thumbs. Write the direct-vs-inverse test on board.
  \item Wed bridge: “Tomorrow we graph these inverse functions and name the features.”
}%
}{%
\phaseitems{
  \item Presenting pair uses sentence frame for Parts A and B.
  \item Rest of class signals and writes takeaway in margin.
}%
}{%
\phaseitems{
  \item Summarize: how do you test direct vs inverse?
  \item Wednesday preview: what shape do you think \(y = k/x\) graphs as?
}%
}{Cold-call a pair that struggled on \#26 — hold them accountable to the direct-vs-inverse distinction.}

\phasetag{summary recap (not graded)}
\frameworkphaseheader{Exit Ticket}{1--2}{3}{%
\phaseitems{
  \item Collect at door.
  \item Scan \#3 (how to distinguish direct from inverse). Plan P3 Do Now if many students are still fuzzy.
}%
}{%
\phaseitems{
  \item 3 sentence stems, honest.
}%
}{%
\phaseitems{
  \item Did students articulate the test (\(y/x\) constant = direct; \(xy\) constant = inverse)?
}%
}{Collect. No grade.}

\begin{callout}{\IconClipboard\ EXIT TICKET — Student Form}{calloutyellow}
Today I learned that \_\_\_.\\
The biggest trap in Ramón's task was \_\_\_.\\
I distinguish direct from inverse variation by \_\_\_.
\end{callout}

\clearpage

\begin{callout}{\IconClock\ PERIOD A / PERIOD F — 65-MIN CLASS NOTES}{calloutpink}
Both sections should have 65 min for P2. Use the extra 10 min on Explore \#26 (Ramón) — the richer discussion, the better the DOK-3 payoff.\\
If finished early: hand out Practice \#25 (SAT/ACT MC) as fast-finisher.
\end{callout}

\begin{callout}{\IconPuzzle\ MODIFICATIONS / SUPPORT FOR IEP STUDENTS}{calloutiep}
\textbullet\ Provide the Application Rules card as a sticker/cut-out.\\
\textbullet\ For \#26, provide pre-formatted work space: “PART A: distance = \_\_\_ \(\times\) time” etc.\\
\textbullet\ Accept verbal sentence-frame completion in lieu of written.
\end{callout}

\begin{callout}{\IconGlobe\ SUPPORTS FOR ELL STUDENTS}{calloutell}
\textbullet\ Sentence frames on student page (don't skip).\\
\textbullet\ Word cards: “frequency” = how often; “volume” = space inside; “pressure” = push.\\
\textbullet\ “Performance Task” = an assessment task with multiple parts.\\
\textbullet\ Pair ELL students with English-dominant partners for TPS.
\end{callout}

\end{document}
